% Section 2.4, Live Demo. Body only: the \subsection{Live Demo} heading stays in the
% main tex. Two variants, variant 1 active, variant 2 commented (same layout as
% 3-1-wikipedia-dataset.tex). The Interactive Use Cases subsubsection that follows
% covers the two workflows (scenarios), so these paragraphs describe interface,
% I/O contract, configurability and implementation only.
%
% Paragraph ladder (important -> details): Functionality (what the user sees and
% does) -> Input and Output (the contract with the user's data) -> Model and
% Prompt Configuration (adaptability, cashes the §2-intro promise "easily
% deployable for new languages and with different backbone models") ->
% Implementation (stack).
%
% Integration steps:
% 1. Add \label{sec:live-demo} to \subsection{Live Demo} (referenced from the
%    2.3 default-criteria sentence).
% 2. Add the tiptap BibTeX entry below to custom.bib.
% 3. Owner decision pending: export formats. Code today exports .xlsx + .md
%    (export.py); the text says "spreadsheet or a Markdown file". If a .txt
%    export ships before submission, swap the sentence ending.
%
% PAPER<->CODE PARITY CHECKLIST (owner committed to implement before the demo
% video / submission; text below already describes the target state):
%   [ ] Refiner config (model + prompt) surface in Settings   — absent today
%   [ ] "Apply all" refinement pass button in the inspector   — today only per-issue accept
%   [ ] Terminology extraction and disambiguation as separately configurable
%       roles — today one shared Grounding card (model + prompt)
%   [ ] (optional) .txt export — today .xlsx + .md
% Facts verified in code 2026-07-10: editable TipTap Translation pane +
% per-paragraph revision history with re-scoring, restore and best-revision
% tracking (EditorParagraph.tsx, app.py /revisions, /restore); upload paste or
% .docx/.md/.txt with enforced equal paragraph counts (file-ingest.ts,
% UploadModal.tsx); AI draft translation job (/documents/{id}/translate);
% model registry name/base_url/api_key/params with OpenRouter + vLLM seed rows
% and a connectivity test (seed.py, app.py /models); per-criterion judge
% config (criterion table); React 19 + TypeScript, custom TipTap extension
% review-extension.ts, FastAPI over single-file SQLite, OpenAI-compatible
% client (llm/client.py), single Docker container.
%
% @misc{tiptap,
%   author       = {{Tiptap}},
%   title        = {Tiptap: the headless rich text editor framework},
%   year         = {2025},
%   howpublished = {\url{https://github.com/ueberdosis/tiptap}},
%   note         = {GitHub repository}
% }

% ===== VARIANT 1 (main) =====
\paragraph{Functionality.}
\systemname{} presents each document as parallel Source and Translation panes
and organizes analysis into four tabs (Document, Glossary, Ranking, Settings).
In the Document tab, findings from the \moduleEvaluation{} module appear as
span-level underlines colored by criterion, each paragraph carries an
aggregate score chip, and an inspector panel lists issues with suggested
corrections. The Translation pane is a directly editable rich-text editor:
users accept suggestions one by one, apply them all in a single refinement
pass, or edit the text manually. Every change produces a new revision that
can be re-scored and restored from the paragraph history. Every extracted
term carries two independent signals, shown as colored spans with hover
popovers: its grounding difficulty on the Source side (the three resolution
states of the \moduleTerminology{} module) and the pair accuracy of its
rendering in the Translation. The Glossary tab groups terms by entity and
exposes their grounding traces, and the Ranking tab orders paragraphs by
score or issue count.

\paragraph{Input and Output.}
A document enters \systemname{} as a pair of Source and Translation aligned
paragraph by paragraph. Users paste text into the editor or upload .docx,
.md, or .txt files for each side, and both sides must contain the same
number of paragraphs. Users can also provide only the Source and generate a
draft with the configured translation model. The resulting translation
exports as a paragraph-aligned spreadsheet or a Markdown file.

\paragraph{Model and Prompt Configuration.}
\systemname{} adapts to new domains through the Settings tab. Its model
registry defines each LLM as an OpenAI-compatible endpoint with a base URL,
an API key, and generation parameters, and ships with preconfigured
OpenRouter and local vLLM entries. Every LLM role in the pipeline is a
reference into this registry paired with an editable prompt: the draft
translator, each judge criterion, the refiner, and the terminology
extraction and disambiguation steps. Defaults tuned for the historical
domain come preloaded, and users can override any of them without
redeployment.

\paragraph{Implementation.}
We build the frontend as a React and TypeScript application. Its editor
builds on the open-source TipTap framework~\cite{tiptap}, which we extend
with a custom decoration layer that draws judge findings as stacked
span-level underlines and marks term spans in place, so all analysis stays
anchored to the text. The backend is a FastAPI service over a single-file
SQLite store and runs translation and first-pass scoring as background jobs.
All LLM calls pass through one OpenAI-compatible client with per-model
endpoint configuration, which makes the system model agnostic: any hosted
gateway, such as OpenRouter, or a local vLLM server can act as the model
backend. Judge scores and issues are stored append-only, so re-evaluation
appends revisions and never overwrites earlier predictions. The system
deploys as a single Docker container with one command.

% ===== VARIANT 2 (text-first framing) =====
% \paragraph{Functionality.}
% \systemname{} concentrates all model output on the text itself. The Document
% tab aligns Source and Translation paragraph by paragraph, underlines each
% detected issue in the color of the criterion that raised it, and shows a
% score chip per paragraph. The Translation pane is not a static view but an
% editable rich-text editor: users review issues in an inspector panel,
% accept suggested corrections one by one or in a single refinement pass, and
% edit the text freely. Each change becomes a revision in the paragraph
% history, where revisions are re-scored against the active criteria and any
% earlier state can be restored. Terms are highlighted with two independent
% signals, grounding difficulty on the Source side and pair accuracy of the
% Translation rendering, and a popover on each term reveals the linked
% Wikidata entity, its resolution state, and the justification. Two further
% tabs support analysis at other granularities: Glossary aggregates terms
% with their grounding traces (query, candidate search, decision steps), and
% Ranking sorts paragraphs by score or issue count for triage.
%
% \paragraph{Input and Output.}
% \systemname{} works on paragraph-aligned document pairs. Both panes accept
% pasted text or .docx, .md, and .txt files, and the upload dialog enforces
% an equal number of paragraphs on the two sides. When only the Source is
% provided, the configured translation model produces the draft Translation.
% At any point the current translation exports as a paragraph-aligned
% spreadsheet or a Markdown file.
%
% \paragraph{Model and Prompt Configuration.}
% Flexibility in \systemname{} comes from one mechanism: every LLM role is a
% model reference plus a prompt, both editable in the Settings tab. The model
% registry declares each LLM as an OpenAI-compatible endpoint (base URL, API
% key, generation parameters) and comes preloaded with OpenRouter and local
% vLLM entries. The draft translator, each judge criterion, the refiner, and
% the terminology extraction and disambiguation steps all draw on this
% registry, so replacing a backbone model or retargeting the system to a new
% domain is a configuration change rather than a code change. Domain-tuned
% defaults are included and remain fully overridable.
%
% \paragraph{Implementation.}
% \systemname{} is model agnostic by construction. Every LLM call, whether
% from the translator, a judge, or the terminology pipeline, goes through one
% OpenAI-compatible client with per-model endpoint configuration, so any
% hosted gateway, such as OpenRouter, or a local vLLM server can serve the
% models. The server is a FastAPI application over a single-file SQLite store
% that executes translation and first-pass scoring as background jobs. Its
% evaluation store is append-only: re-evaluation appends new score revisions
% and never overwrites earlier predictions. The client is a React and
% TypeScript application whose editor builds on the open-source TipTap
% framework~\cite{tiptap}, extended with a custom decoration layer that
% stacks criterion-colored underlines and term marks over the Translation
% text. The full system ships as one Docker container and starts with a
% single command.
